% (c) Nikita Lisitsa, lisyarus@gmail.com, 2026

\documentclass[10pt]{beamer}

\usepackage[T2A]{fontenc}
\usepackage[russian]{babel}
\usepackage{minted}

\usepackage[most]{tcolorbox}
\tcbuselibrary{minted}

\usepackage{graphicx}
\graphicspath{ {./images/} }

\usepackage{tikz}
\usepackage{xifthen}

\usepackage{adjustbox}

\usepackage{color}
\usepackage{soul}

\usepackage{hyperref}

\usetheme{metropolis}

\definecolor{red}{rgb}{1,0,0}
\definecolor{green}{rgb}{0,0.5,0}
\definecolor{blue}{rgb}{0,0,1}
\definecolor{codebg}{RGB}{29,35,49}
\setminted{fontsize=\footnotesize}

\usemintedstyle{solarized-light}

\makeatletter
\newcommand{\slideimage}[1]{
  \begin{figure}
    \begin{adjustbox}{width=\textwidth, totalheight=\textheight-2\baselineskip-2\baselineskip,keepaspectratio}
      \includegraphics{#1}
    \end{adjustbox}
  \end{figure}
}
\makeatother

\title{Компьютерная графика}
\subtitle{Лекция 3: Буферы на GPU, аттрибуты вершин, индексированный рендеринг, instancing}
\date{2026}

\setbeamertemplate{footline}[frame number]

\begin{document}

\frame{\titlepage}

\begin{frame}[fragile]
\frametitle{Вершины}
\begin{itemize}
\item Мы знаем, что \textit{вершинный шейдер} обрабатывает по одной \textit{вершины} нашей 3D (или 2D) модели
\pause
\item Пока мы умеем только хардкодить их \textit{атрибуты} (положение в пространстве, цвет) в шейдере
\pause
\item В модели могут быть тысячи или миллионы вершин, одним шейдером хочется рисовать разные модели
\pause
\item \(\Longrightarrow\) Хочется как-то загрузить их в память GPU
\pause
\item \(\Longrightarrow\) Хочется как-то рассказать шейдеру, где они лежат
\end{itemize}
\end{frame}

\begin{frame}[fragile]
\frametitle{Буферы}
\usemintedstyle{solarized-light}
\begin{itemize}
\item Могут хранить произвольные данные на GPU, сами никак их не интерпретируют
\pause
\item При создании нужно указать:
\pause
\begin{itemize}
\item \mintinline{rust}|size| -- размер в байтах
\pause
\item \mintinline{rust}|usage| -- флаги, описывающие, как мы будем использовать этот буфер (нужны для валидации, выбора правильного вида памяти, и т.д.)
\end{itemize}
\end{itemize}
\end{frame}

\begin{frame}[fragile]
\frametitle{Buffer usage}
\usemintedstyle{solarized-light}
\begin{itemize}
\item Возможные значения \mintinline{rust}|usage| (это битовые флаги -- их можно комбинировать):
\pause
\begin{itemize}
\item \mintinline{rust}|MapRead|, \mintinline{rust}|MapWrite| -- можно замапить на чтение/запись (об этом позже)
\pause
\item \mintinline{rust}|CopySrc|, \mintinline{rust}|CopyDst| -- можно копировать данные из/в этот буффер
\pause
\item \mintinline{rust}|Index|, \mintinline{rust}|Vertex| -- можно использовать для индексов / самих вершин
\pause
\item \mintinline{rust}|Uniform| -- можно использовать для \textit{uniform значений} (как immediates, но больше объём)
\pause
\item \mintinline{rust}|Storage| -- можно использовать для \textit{random access storage}
\pause
\item \mintinline{rust}|Indirect| -- можно использовать для \textit{indirect draw} команд
\pause
\item \mintinline{rust}|QueryResolve| -- можно использовать для сохранения результатов \textit{occlusion/timestamp queries}
\end{itemize}
\end{itemize}
\end{frame}

\begin{frame}[fragile]
\frametitle{Buffer size}
\usemintedstyle{solarized-light}
\begin{itemize}
\item Создание буфера работает примерно как \mintinline{cpp}|malloc|: память выделяется один раз, перевыделить её нельзя
\pause
\item Если вам нужен буфер другого размера, придётся его пересоздать
\pause
\item Для динамических данных (создаются/удаляются объекты и т.п.) придётся руками реализовывать обновление размера (примерно как в \mintinline{cpp}|std::vector|)
\end{itemize}
\end{frame}

\begin{frame}[fragile]
\frametitle{Буферы: запись данных}
\begin{itemize}
\item Записать данные в буфер можно через \mintinline{cpp}|queue|:
\usemintedstyle{lightbulb}
\begin{tcblisting}{listing engine=minted, minted language=cpp, minted options={fontsize=\scriptsize}, listing only, colback=codebg, colframe=codebg, rounded corners}
void wgpuQueueWriteBuffer(WGPUQueue queue, WGPUBuffer buffer,
    uint64_t bufferOffset, void const * data, size_t size);
\end{tcblisting}
\usemintedstyle{solarized-light}
\begin{itemize}
\item \mintinline{cpp}|buffer| -- в какой буфер грузим данные
\item \mintinline{cpp}|bufferOffset| -- начиная с какого байта в буфере загрузить данные
\item \mintinline{cpp}|data| -- указатель на данные
\item \mintinline{cpp}|size| -- размер данных в байтах
\end{itemize}
\end{itemize}
\end{frame}

\begin{frame}[fragile]
\frametitle{Буферы: запись данных}
\begin{itemize}
\item Так же, как с другими операциями, \mintinline{rust}|queue.writeBuffer| не пишет сразу в память GPU, а кладёт в очередь команду записи
\pause
\item После вызова \mintinline{rust}|queue.writeBuffer| с данными по указателю \mintinline{cpp}|data| можно делать всё, что угодно (в т.ч. удалить)
\pause
\item \begin{math}\Longrightarrow\end{math} WebGPU сначала копирует данные в собственную память, а потом инициирует отложенную загрузку на GPU
\end{itemize}
\end{frame}

\begin{frame}[fragile]
\frametitle{Буферы: запись данных}
\begin{itemize}
\item Если мы обновляем данные в буфере каждый кадр, GPU придётся ждать, пока буфер перестанет быть нужен, чтобы обновить в нём данные
\pause
\item \begin{math}\Longrightarrow\end{math} WebGPU, скорее всего, внутри будет делать двойную (или тройную) буферизацию -- иметь несколько копий буфера, циклически переиспользующихся от кадра к кадру
\end{itemize}
\end{frame}

\begin{frame}[fragile]
\frametitle{Mapped buffer}
\usemintedstyle{solarized-light}
\begin{itemize}
\item Можно получить виртуальный указатель на буфер или его часть и пользоваться им для чтения/записи: \mintinline{rust}|buffer.mapAsync(mode, offset, size, callback)|
\pause
\begin{itemize}
\item \mintinline{cpp}|mode| -- комбинация флагов \mintinline{cpp}|READ| и \mintinline{cpp}|WRITE|
\pause
\item \mintinline{cpp}|offset|, \mintinline{cpp}|size| -- какой диапазон буфера хотим замапить
\pause
\item \mintinline{cpp}|callback| -- вызывается, когда буфер будет замаплен (операция \textit{асинхронная})
\end{itemize}
\pause
\item Когда буфер запамлен:
\begin{itemize}
\item \mintinline{cpp}|buffer.getMappedRange/get_mapped_range_mut| -- получить указатель на чтение/запись
\item \mintinline{cpp}|buffer.getConstMappedRange/get_mapped_range| -- получить указатель на чтение
\end{itemize}
\pause
\item После использования буфер нужно отмапить: \mintinline{cpp}|buffer.unmap()|
\pause
\item Пока буфер замаплен, работать с буфером (загружать данные другими методами, использовать данные для рисования) нельзя
\end{itemize}
\end{frame}

\begin{frame}[fragile]
\frametitle{Mapped buffer}
\usemintedstyle{solarized-light}
\begin{itemize}
\item Mapped buffers работают через UMA/ReBAR, где это поддерживается, либо через явное копирование CPU \(\leftrightarrow\) GPU
\pause
\item Единственный способ \textit{прочитать} что-то из GPU буфера на CPU
\pause
\item Буфер можно создать с флагом \mintinline{cpp}|mappedAtCreation|, тогда он будет сразу замаплен (без асинхронного ожидания), -- удобно для буферов, данные которых загружаются только один раз при создании
\end{itemize}
\end{frame}

\begin{frame}[fragile]
\frametitle{Буферы: типичный пример создания}
\usemintedstyle{lightbulb}
\begin{tcblisting}{listing engine=minted, minted language=rust, minted options={fontsize=\scriptsize}, listing only, colback=codebg, colframe=codebg, rounded corners}
// Псевдокод :)

vertex_buffer = device.create_buffer(
    size = vertices.size() * sizeof(vertex),
    usage = VERTEX | COPY_DST
);

queue.write_buffer(&vertex_buffer, 0, vertex_buffer.size());
\end{tcblisting}
\end{frame}

\begin{frame}[fragile]
\frametitle{Атрибуты вершин}
\begin{itemize}
\item Вершины -- основные входные данные графического конвейера
\pause
\item С точки зрения графических API, вершины характеризуются набором \textit{атрибутов}: позиция, цвет, нормаль, текстурные координаты, и т.п.
\pause
\item Каждый атрибут описывается отдельно, может иметь свой тип и формат хранения, -- всё это запоминается в \textit{render pipeline}
\end{itemize}
\end{frame}

\begin{frame}[fragile]
\frametitle{Атрибуты вершин: шейдер}
\usemintedstyle{lightbulb}
\begin{tcblisting}{listing engine=minted, minted language=wgsl, minted options={fontsize=\scriptsize}, listing only, colback=codebg, colframe=codebg, rounded corners}
struct VertexIn
{
    @location(0) position : vec3f,
    @location(1) normal : vec3f,
    @location(2) color : vec4f,
    @location(3) texcoord : vec2f,
}
\end{tcblisting}
\pause
\usemintedstyle{solarized-light}
\begin{itemize}
\item Location: \mintinline{cpp}|0, 1, ... maxVertexAttributes - 1|
\pause
\item \mintinline{cpp}|maxVertexAttributes| \begin{math}\geq\end{math} 16
\pause
\item Есть тип (\mintinline{rust}|f32|/\mintinline{rust}|i32|/\mintinline{rust}|u32|) и размерность (\mintinline{rust}|1|/\mintinline{rust}|2|/\mintinline{rust}|3|/\mintinline{rust}|4|)
\end{itemize}
\end{frame}

\begin{frame}[fragile]
\frametitle{Атрибуты вершин: формат}
\usemintedstyle{solarized-light}
Floating-point (\mintinline{rust}|f32|/\mintinline{rust}|vec2f|/...) форматы:
\begin{center}
\begin{tabular}{ |c|c|c| } 
\hline
 & 16-bit & 32-bit \\ 
\hline
1D & \mintinline{rust}|Float16| & \mintinline{rust}|Float32| \\ 
\hline
2D & \mintinline{rust}|Float16x2| & \mintinline{rust}|Float32x2| \\ 
\hline
3D & & \mintinline{rust}|Float32x3| \\ 
\hline
4D & \mintinline{rust}|Float16x4| & \mintinline{rust}|Float32x4| \\ 
\hline
\end{tabular}
\end{center}
\end{frame}

\begin{frame}[fragile]
\frametitle{Атрибуты вершин: формат}
\usemintedstyle{solarized-light}
Unsigned integer (\mintinline{rust}|u32|/\mintinline{rust}|vec2u|/...) форматы:
\begin{center}
\begin{tabular}{ |c|c|c|c| } 
\hline
 & 8-bit & 16-bit & 32-bit \\ 
\hline
1D & \mintinline{rust}|Uint8| & \mintinline{rust}|Uint16| & \mintinline{rust}|Uint32| \\ 
\hline
2D & \mintinline{rust}|Uint8x2| & \mintinline{rust}|Uint16x2| & \mintinline{rust}|Uint32x2| \\ 
\hline
3D & & & \mintinline{rust}|Uint32x3| \\ 
\hline
4D & \mintinline{rust}|Uint8x4| & \mintinline{rust}|Uint16x4| & \mintinline{rust}|Uint32x4| \\ 
\hline
\end{tabular}
\end{center}
\end{frame}

\begin{frame}[fragile]
\frametitle{Атрибуты вершин: формат}
\usemintedstyle{solarized-light}
Signed integer (\mintinline{rust}|i32|/\mintinline{rust}|vec2i|/...) форматы:
\begin{center}
\begin{tabular}{ |c|c|c|c| } 
\hline
 & 8-bit & 16-bit & 32-bit \\ 
\hline
1D & \mintinline{rust}|Sint8| & \mintinline{rust}|Sint16| & \mintinline{rust}|Sint32| \\ 
\hline
2D & \mintinline{rust}|Sint8x2| & \mintinline{rust}|Sint16x2| & \mintinline{rust}|Sint32x2| \\ 
\hline
3D & & & \mintinline{rust}|Sint32x3| \\ 
\hline
4D & \mintinline{rust}|Sint8x4| & \mintinline{rust}|Sint16x4| & \mintinline{rust}|Sint32x4| \\ 
\hline
\end{tabular}
\end{center}
\end{frame}

\begin{frame}[fragile]
\frametitle{Атрибуты вершин: нормированные форматы}
\usemintedstyle{solarized-light}
\begin{itemize}
\item Floating-point форматы часто слишком расточительны: большой диапазон, относительная точность
\pause
\item Например, цвет часто хранится 8-битным числом 0..255 и интерпретируется как число \begin{math}[0, 1]\end{math} (по формуле \mintinline{cpp}|x / 255.0|)
\pause
\item Такие форматы называются \textit{нормированными}: хранятся как знаковое/беззнаковое целое, но в шейдер попадают как floating-point (преобразование за нас делает GPU)
\pause
\begin{itemize}
\item Для беззнаковых: \begin{math}[0, 2^{bits}-1] \rightarrow [0, 1]\end{math} по формуле \mintinline{cpp}|x / pow(2, bits)|
\pause
\item Для знаковых: \begin{math}[-2^{bits-1}, 2^{bits-1}-1] \rightarrow [-1, 1]\end{math} по формуле \mintinline{cpp}|clamp(x / (pow(2, bits-1) - 1), -1, 1)|
\end{itemize}
\pause
\item Полезно для компактного хранения атрибутов (цвет, нормали)
\pause
\item Имеет смысл \textit{только} для floating-point аттрибутов
\end{itemize}
\end{frame}

\begin{frame}[fragile]
\frametitle{Атрибуты вершин: нормированные форматы}
\usemintedstyle{solarized-light}
Unsigned integer нормированные форматы:
\begin{center}
\begin{tabular}{ |c|c|c|c| } 
\hline
 & 8-bit & 16-bit & 32-bit \\ 
\hline
1D & \mintinline{rust}|Unorm8| & \mintinline{rust}|Unorm16| & \mintinline{rust}|Unorm32| \\ 
\hline
2D & \mintinline{rust}|Unorm8x2| & \mintinline{rust}|Unorm16x2| & \mintinline{rust}|Unorm32x2| \\ 
\hline
3D & & & \mintinline{rust}|Unorm32x3| \\ 
\hline
4D & \mintinline{rust}|Unorm8x4| & \mintinline{rust}|Unorm16x4| & \mintinline{rust}|Unorm32x4| \\ 
\hline
\end{tabular}
\end{center}
\end{frame}

\begin{frame}[fragile]
\frametitle{Атрибуты вершин: нормированные форматы}
\usemintedstyle{solarized-light}
Signed integer нормированные форматы:
\begin{center}
\begin{tabular}{ |c|c|c|c| } 
\hline
 & 8-bit & 16-bit & 32-bit \\ 
\hline
1D & \mintinline{rust}|Snorm8| & \mintinline{rust}|Snorm16| & \mintinline{rust}|Snorm32| \\ 
\hline
2D & \mintinline{rust}|Snorm8x2| & \mintinline{rust}|Snorm16x2| & \mintinline{rust}|Snorm32x2| \\ 
\hline
3D & & & \mintinline{rust}|Snorm32x3| \\ 
\hline
4D & \mintinline{rust}|Snorm8x4| & \mintinline{rust}|Snorm16x4| & \mintinline{rust}|Snorm32x4| \\ 
\hline
\end{tabular}
\end{center}
\end{frame}

\begin{frame}[fragile]
\frametitle{Атрибуты вершин: буферы}
\usemintedstyle{solarized-light}
\begin{itemize}
\item Разные атрибуты могут лежать в разных буферах, поэтому в \mintinline{cpp}|renderPipelineDescriptor.vertex| описывается массив \textit{buffer layouts}
\pause
\item Каждый \textit{buffer layout} описывает массив атрибутов, лежащих вместе в одном буфере
\pause
\item У каждого атрибута есть сдвиг (\textit{offset}) от начала буфера до значения атрибута нулевой вершины
\pause
\item У всех атрибутов одного буфера есть общее расстояние (\textit{stride}) между значениями атрибутов соседних вершин
\end{itemize}
\end{frame}

\begin{frame}[fragile]
\frametitle{Атрибуты вершин: stride}
\usemintedstyle{solarized-light}
\begin{itemize}
\item Аттрибут вершины с номером \mintinline{cpp}|i| будет взят по смещению \mintinline{cpp}|offset + stride * i|
\pause
\begin{itemize}
\item Нулевая вершина: \mintinline{cpp}|offset|
\item Первая вершина: \mintinline{cpp}|offset + stride|
\item Вторая вершина: \mintinline{cpp}|offset + 2 * stride|
\item И т.д.
\end{itemize}
\pause
\item Нужен для гибкости хранения аттрибутов
\end{itemize}
\end{frame}

\begin{frame}[fragile]
\frametitle{Атрибуты вершин: пример array-of-structs, один буфер}
\usemintedstyle{lightbulb}
\begin{tcblisting}{listing engine=minted, minted language=cpp, minted options={fontsize=\scriptsize}, listing only, colback=codebg, colframe=codebg, rounded corners}
struct vertex {
    float position[3];
    float normal[3];
    unsigned char color[4];
};

vertex_buffers = [
    {
        stride: sizeof(vertex),
        attributes: [
            { format: Float32x3, offset:  0, location: 0 },
            { format: Float32x3, offset: 12, location: 1 },
            { format:  Unorm8x4, offset: 24, location: 2 },
        ],
    },
];
\end{tcblisting}
\begin{center}
\begin{tikzpicture}
\foreach \x
[evaluate=\x as \ta using int(\x * 28)]
[evaluate=\x as \tb using int(\x * 28 + 12)]
[evaluate=\x as \tc using int(\x * 28 + 24)]
[evaluate=\x as \td using int(\x * 28 + 28)]
in {0, 1, 2}
{
\filldraw[fill=red!40!white, draw=black,thick] (3.0 * \x + 0.0, 0) rectangle (3.0 * \x + 1.0, 0.5);
\filldraw[fill=green!40!white, draw=black,thick] (3.0 * \x + 1.0, 0) rectangle (3.0 * \x + 2.0, 0.5);
\filldraw[fill=blue!40!white, draw=black,thick] (3.0 * \x + 2.0, 0) rectangle (3.0 * \x + 3.0, 0.5);
\node[] at (3.0 * \x + 0.5, 0.25) {P\x};
\node[] at (3.0 * \x + 1.5, 0.25) {N\x};
\node[] at (3.0 * \x + 2.5, 0.25) {C\x};
\node[] at (3.0 * \x + 0.0, 0.75) {\ta};
\node[] at (3.0 * \x + 1.0, 0.75) {\tb};
\node[] at (3.0 * \x + 2.0, 0.75) {\tc};
\node[] at (3.0 * \x + 3.0, 0.75) {\td};
\draw[black,thick] (3.0 * \x, -0.125) -- (3.0 * \x, -0.5);
\draw[black,thick] (3.0 * \x + 3.0, -0.125) -- (3.0 * \x + 3.0, -0.5);
\node[] at (3.0 * \x + 1.5, -0.25) {v\x};
}
\end{tikzpicture}
\end{center}
\end{frame}

\begin{frame}[fragile]
\frametitle{Аттрибуты вершин: пример struct-of-arrays, три буфера}
\usemintedstyle{lightbulb}
\begin{tcblisting}{listing engine=minted, minted language=cpp, minted options={fontsize=\scriptsize}, listing only, colback=codebg, colframe=codebg, rounded corners}
float positions[3 * N];
float normals[3 * N];
unsigned char colors[4 * N];

vertex_buffers = [
    {
        stride: 12,
        attributes: [ { format: Float32x3, offset: 0, location: 0 } ],
    },
    {
        stride: 12,
        attributes: [ { format: Float32x3, offset: 0, location: 1 } ],
    },
    {
        stride: 4,
        attributes: [ { format:  Unorm8x4, offset: 0, location: 2 } ],
    },
];
\end{tcblisting}
\begin{center}
\begin{tikzpicture}
\node[] at (0.0, 0.75) {0};
\node[] at (4.0, 0.75) {0};
\node[] at (8.0, 0.75) {0};
\foreach \x
[evaluate=\x as \ta using int(\x * 12 + 12)]
[evaluate=\x as \tb using int(\x * 12 + 12)]
[evaluate=\x as \tc using int(\x * 4 + 4)]
in {0, 1, 2}
{
\filldraw[fill=red!40!white, draw=black,thick] (\x + 0.0, 0) rectangle (\x + 1.0, 0.5);
\filldraw[fill=green!40!white, draw=black,thick] (\x + 4.0, 0) rectangle (\x + 5.0, 0.5);
\filldraw[fill=blue!40!white, draw=black,thick] (\x + 8.0, 0) rectangle (\x + 9.0, 0.5);
\node[] at (\x + 0.5, 0.25) {P\x};
\node[] at (\x + 4.5, 0.25) {N\x};
\node[] at (\x + 8.5, 0.25) {C\x};
\node[] at (\x + 1.0, 0.75) {\ta};
\node[] at (\x + 5.0, 0.75) {\tb};
\node[] at (\x + 9.0, 0.75) {\tc};
}
\end{tikzpicture}
\end{center}
\end{frame}

\begin{frame}[fragile]
\frametitle{Аттрибуты вершин: как хранить?}
\usemintedstyle{solarized-light}
\begin{itemize}
\item По смыслу, каждый \textit{buffer layout} описывает массив структур (\textit{stride} -- размер структуры)
\pause
\item По умолчанию используйте один буфер со структурой вершины (например, если вершины хранятся непрерывано в памяти, e.g. \mintinline{cpp}|std::vector<vertex>|)
\pause
\item Если часть атрибутов постоянны, а другую часть нужно иногда обновлять -- разделите эти части по двум разным буферам
\pause
\item Если часть атрибутов не нужны для некоторых алгоритмов (например, для теней) -- выделите их в отдельный буфер
\pause
\item Некоторые форматы 3D сцен (напр. \href{https://registry.khronos.org/glTF/specs/2.0/glTF-2.0.html}{\texttt{glTF}}) хранят вершины в бинарном формате, удобном для загрузки на GPU: поддерживают несколько буферов, разный \mintinline{cpp}|stride|, и разные форматы атрибутов
\end{itemize}
\end{frame}

\begin{frame}[fragile]
\frametitle{Буферы с вершинами}
\usemintedstyle{solarized-light}
\begin{itemize}
\item Для рендеринга нужно выставить буфер с вершинами через \mintinline{rust}|render_pass.set_vertex_buffer(slot, buffer, offset, size)|
\pause
\item \mintinline{rust}|slot| -- номер \textit{buffer layout}'а в массиве \mintinline{cpp}|renderPipelineDescriptor.vertex|
\pause
\item \mintinline{rust}|offset|, \mintinline{rust}|size| -- позволяют использовать только часть буфера (offset'ы атрибутов будут браться относительно начала этой части)
\end{itemize}
\end{frame}

\begin{frame}[fragile]
\frametitle{Аттрибуты вершин: пример полностью}
\usemintedstyle{lightbulb}
\begin{tcblisting}{listing engine=minted, minted language=cpp, minted options={fontsize=\scriptsize}, listing only, colback=codebg, colframe=codebg, rounded corners}
// Инициализация:
pipeline = createPipeline()
vertexBuffer = loadBuffer(vertices, Usage::Vertex)

// Рендеринг:
renderPass.setPipeline(pipeline)
renderPass.setVertexBuffer(vertexBuffer)
renderPass.draw(vertices.size())
\end{tcblisting}
\end{frame}

\begin{frame}[fragile]
\frametitle{Индексированный рендеринг}
\usemintedstyle{solarized-light}
\slideimage{spheres.png}
\begin{itemize}
\item В моделях одна вершина часто входит в несколько треугольников (в среднем, 6)
\pause
\item Если использовать примитив \mintinline{cpp}|TriangleList|, придётся копировать каждую вершину несколько раз
\pause
\item \mintinline{cpp}|TriangleStrip| только частично решают проблему
\pause
\item \begin{math}\Longrightarrow\end{math} Индексированный рендеринг
\end{itemize}
\end{frame}

\begin{frame}[fragile]
\frametitle{Индексированный рендеринг}
\usemintedstyle{solarized-light}
\begin{itemize}
\item \mintinline{cpp}|renderPass.draw(first=3, count=5)|:
\begin{center}
\begin{tikzpicture}
\foreach \x in {0, ..., 9}
{
\ifthenelse{\x > 2 \AND \x < 8}
{\filldraw[fill=red!40!white, draw=black,thick] (\x, 0) rectangle (\x + 1.0, 0.5);}
{\draw[black,thick] (\x, 0) rectangle (\x + 1.0, 0.5);}

\node[] at (\x + 0.5, 0.25) {v\x};
}
\end{tikzpicture}
\end{center}
\pause

\vspace{1cm}

\item \mintinline{cpp}|renderPass.draw_indexed(first=3, count=5)|:
\begin{center}
\begin{tikzpicture}
\foreach \x in {0, ..., 9}
{
\ifthenelse{\x = 6 \OR \x = 3 \OR \x = 4 \OR \x = 0 \OR \x = 8}
{\filldraw[fill=red!40!white, draw=black,thick] (\x, 2) rectangle (\x + 1.0, 2.5);}
{\draw[black,thick] (\x, 2) rectangle (\x + 1.0, 2.5);}
\node[] at (\x + 0.5, 2.25) {v\x};
}

\foreach \x in {0, ..., 9}
{
\ifthenelse{\x > 2 \AND \x < 8}
{\filldraw[fill=red!40!white, draw=black,thick] (\x, 0) rectangle (\x + 1.0, 0.5);}
{\draw[black,thick] (\x, 0) rectangle (\x + 1.0, 0.5);}
}
\node[] at (0.5, 0.25) {3};
\node[] at (1.5, 0.25) {0};
\node[] at (2.5, 0.25) {5};
\node[] at (3.5, 0.25) {6};
\node[] at (4.5, 0.25) {3};
\node[] at (5.5, 0.25) {4};
\node[] at (6.5, 0.25) {0};
\node[] at (7.5, 0.25) {8};
\node[] at (8.5, 0.25) {0};
\node[] at (9.5, 0.25) {1};

\draw[thick,-stealth] (3.5, 0.5) -- (6.5, 1.9);
\draw[thick,-stealth] (4.5, 0.5) -- (3.5, 1.9);
\draw[thick,-stealth] (5.5, 0.5) -- (4.5, 1.9);
\draw[thick,-stealth] (6.5, 0.5) -- (0.5, 1.9);
\draw[thick,-stealth] (7.5, 0.5) -- (8.5, 1.9);

% 0 1 2 3 4 5 6 7 8 9
% 3 0 5[6 3 4 0 8]0 1
\end{tikzpicture}
\end{center}
\end{itemize}
\end{frame}

\begin{frame}[fragile]
\frametitle{Индексированный рендеринг: квадрат}
\usemintedstyle{solarized-light}
\begin{center}
\begin{tikzpicture}
\draw[black,thick] (0, 0) rectangle (3, 3);
\draw[black,thick] (0, 3) -- (3, 0);
\node[] at (-0.25, -0.25) {0};
\node[] at (3.25, -0.25) {1};
\node[] at (-0.25, 3.25) {2};
\node[] at (3.25, 3.25) {3};
\end{tikzpicture}
\end{center}
\pause
\begin{itemize}
\item Без индексов, \mintinline{cpp}|TriangleList|:
\begin{itemize}
\item Вершины: \mintinline{cpp}|v0, v1, v2, v2, v1, v3|
\end{itemize}
\pause
\item Без индексов, \mintinline{cpp}|TriangleStrip|:
\begin{itemize}
\item Вершины: \mintinline{cpp}|v0, v1, v2, v3|
\end{itemize}
\pause
\item С индексами, \mintinline{cpp}|TriangleList|:
\begin{itemize}
\item Вершины: \mintinline{cpp}|v0, v1, v2, v3|
\item Индексы: \mintinline{cpp}|0, 1, 2, 2, 1, 3|
\end{itemize}
\end{itemize}
\end{frame}

\begin{frame}[fragile]
\frametitle{Индексированный рендеринг}
\usemintedstyle{solarized-light}
\begin{itemize}
\item Индексы хранятся в буфере с \mintinline{cpp}|usage = Index|
\pause
\item \mintinline{rust}|render_pass.set_index_buffer(buffer, format, offset, size)| устанавливает текущий буфер с индексами
\pause
\item Форматов только два: \mintinline{rust}|Uint16| и \mintinline{rust}|Uint32|
\pause
\item \mintinline{rust}|render_pass.draw_indexed()| вызывает индексированный рендеринг
\end{itemize}
\end{frame}

\begin{frame}[fragile]
\frametitle{Индексированный рендеринг: пример}
\usemintedstyle{lightbulb}
\begin{tcblisting}{listing engine=minted, minted language=cpp, minted options={fontsize=\scriptsize}, listing only, colback=codebg, colframe=codebg, rounded corners}
// Инициализация:
pipeline = createPipeline()
vertexBuffer = loadBuffer(vertices, Usage::Vertex)
indexBuffer = loadBuffer(indices, Usage::Index)

// Рендеринг:
renderPass.setPipeline(pipeline)
renderPass.setVertexBuffer(vertexBuffer)
renderPass.setIndexBuffer(indexBuffer)
renderPass.drawIndexed(indices.size())
\end{tcblisting}
\end{frame}

\begin{frame}[fragile]
\frametitle{Primitive restart}
\usemintedstyle{solarized-light}
\begin{itemize}
\item При использовании \mintinline{cpp}|LineStrip| или \mintinline{cpp}|TriangleList| все входные вершины образуют один line strip/triangle strip
\pause
\item Что, если мы хотим одну последовательность вершин разбить на несколько примитивов?
\pause
\item Можно сделать несколько вызовов \mintinline{cpp}|render_pass.draw(...)|, но это расточительно
\pause
\item Максимально возможное для этого типа (\mintinline{cpp}|uint16|/\mintinline{cpp}|uint32|) значение индекса означает что со следующего индекса нужно начать новый примитив:
\mintinline{python}|[0, 1, 2, 3, 65535, 4, 5, 6, 7]| эквивалентно
\mintinline{python}|[0, 1, 2, 3], [4, 5, 6, 7]|
\end{itemize}
\end{frame}

\begin{frame}[fragile]
\frametitle{Рендеринг нескольких объектов}
\begin{itemize}
\item Для каждого \textit{типа объекта} (одинаковый набор атрибутов, шейдеры) -- свой pipeline
\pause
\item Для каждого \textit{отдельного объекта} -- свой набор буферов (напр. vertex + index)
\pause
\item Можно иметь общий буфер на все вершины и индексы всех объектов, если их не нужно обновлять, и делать \mintinline{rust}|render_pass.set_vertex_buffer(0, buffer, offset=...)|
\end{itemize}
\end{frame}

\begin{frame}[fragile]
\frametitle{Instancing}
\begin{itemize}
\item У настройки \textit{buffer layout}'а есть ещё один параметр -- \textit{step mode}
\pause
\item Значение \mintinline{cpp}|stepMode = vertex| означает, что атрибуты будут браться по \textit{номеру вершины}
\pause
\item Значение \mintinline{cpp}|stepMode = instance| означает, что атрибуты будут браться по \textit{номеру инстанса}
\end{itemize}
\end{frame}

\begin{frame}[fragile]
\frametitle{Instancing}
\begin{itemize}
\item \textit{Instancing}, он же \textit{instanced rendering} позволяет одним вызовом нарисовать много копий объекта
\pause
\item Типичная разбивка атрибутов:
\begin{itemize}
\item \mintinline{cpp}|stepMode = vertex| -- атрибуты вершин 3D-модели: позиции, нормали, цвета
\item \mintinline{cpp}|stepMode = instance| -- атрибуты каждой отдельной копии: позиция, поворот, дополнительный цвет
\end{itemize}
\pause
\item Как правило, vertex и instance атрибуты лежат в разных буферах
\pause
\item Вызов \mintinline{cpp}|render_pass.draw(vertex_count=N, instance_count=M)| обработает \(N\cdot M\) вершин: \(M\) копий модели из \(N\) вершин
\end{itemize}
\end{frame}

\begin{frame}[fragile]
\frametitle{Instancing}
\begin{itemize}
\item Раньше instancing имел значительные накладные расходы, и использовался только для большого количества копий (сотен или тысяч)
\pause
\item В современных API instancing -- основная функциональность, все функции рисования (\mintinline{cpp}|render_pass.draw()|) по умолчанию работают с instancing'ом
\end{itemize}
\end{frame}

\end{document}
